\subsubsection{Distribution stationnaire}

La distribution stationnaire d'une chaîne de Markov, 
notée $\pi_\infty$, est une distribution de probabilité qui 
reste invariante lors de l'application de la matrice de transition $P$. 
Mathématiquement, elle satisfait le système d'équations suivant :
\begin{equation}
    \pi_\infty P = \pi_\infty \quad \text{soumis à la contrainte} \quad \sum_{i \in \mathcal{S}} (\pi_\infty)_i = 1
\end{equation}

D'un point de vue algébrique, cela signifie que $\pi_\infty$ est le vecteur 
propre à gauche de la matrice $P$ associé à la valeur propre $\lambda = 1$.

Plutôt que d'estimer cette distribution en élevant $P$ à une grande puissance 
(comme fait précédemment), nous l'avons calculée analytiquement en procédant 
à la décomposition spectrale de la transposée de la matrice $P$ ($P^T$).

Après extraction du vecteur propre correspondant à la valeur 
propre $\lambda = 1$ et normalisation pour que la somme de ses éléments vaille 1, 
nous obtenons la distribution stationnaire suivante :
\begin{equation}
    \pi_\infty = 
    \begin{pmatrix}
        \pi_A \\ \pi_C \\ \pi_G \\ \pi_T
    \end{pmatrix}^T
    = 
    \begin{pmatrix}
        0.2881 \\ 0.3497 \\ 0.2091 \\ 0.1531 
    \end{pmatrix}^T
\end{equation}

\subparagraph{Conclusion sur la convergence}
Nous constatons que ce calcul analytique direct corrobore parfaitement 
l'observation empirique de la question précédente (1.1.2). Les valeurs de $\pi_\infty$ 
correspondent exactement aux probabilités vers lesquelles les lignes de la matrice $P^{50}$ 
avaient convergé. La chaîne de Markov modélisant notre séquence d'ADN possède donc 
bien une unique distribution stationnaire vers laquelle le système tend asymptotiquement, 
quelle que soit sa condition initiale.